Pretraining paradigm is widely assumed to amplify spurious correlations in contrastive self-supervised learning, where augmentation invariance makes salient spurious features --- like background texture --- strongly instance-discriminative. We test this assumption directly with a controlled 4-paradigm comparison on frozen ResNet-50 representations across two group-annotated benchmarks. Contrary to theory, supervised ERM encodes spurious background features \textit{more} strongly than contrastive MoCo-v3 on frozen ResNet-50 representations (Cohen's $d = 5.68$, ANOVA $F = 35.99$, $p = 2.42 \times 10^{-7}$), while self-distillation (DINO) matches ERM despite using no explicit class labels. The key insight is that supervised label correlation --- not augmentation-based instance-discriminability --- is the dominant driver of spurious encoding: label-agnostic pretraining partially suppresses spurious features as a free byproduct of not seeing class labels. On CelebA, the paradigm ranking reverses (MoCo-v3 highest, DINO lowest), revealing a spurious-attribute-type $\times$ pretraining-objective interaction that requires a cross-dataset design to detect. Our findings provide practitioners with evidence-based backbone selection guidance: contrastive SSL backbones partially suppress spurious features by design, without any robustness intervention.
